\section{Resources and frameworks}\label{sec:resources}


\regime{Control}{what each application domain demands of persistent state, read as a stress test.}
Persistent state, its lifecycle, and its governance look different once a domain supplies the substrate on which an always-on agent actually runs. The persistence problem does not arrive in a generic form. A clinical agent and a web agent both accumulate, retrieve, and act on state, but answer to different obligations: the clinical agent must reconcile a patient's recall-biased self-report against an authoritative record and survive a regulator's audit; the web agent must not let a stale workflow abstraction fire an irreversible purchase. Each domain therefore supplies a stress test that sharpens one or more of the six axes (authority, scope, mutability, provenance, recoverability, actionability) and one or more of the five invariants (authority monotonicity, scope non-expansion, deletion propagation, provenance preservation, rollback traceability). The question is not what an agent in the domain does, but what persistence demand the domain adds, which benchmarks and systems expose it, and where the governance gap remains open.

The domain view shows that the skew is not a property of the methods literature alone. Applications inherit it. Domains race to build agents that remember more, recall faster, and personalize more aggressively, and they build benchmarks that score recall and task success. The return arc of the lifecycle, which validates a write before it lands, propagates a deletion, preserves provenance through consolidation, and traces a rollback, is thin in every domain we examine, and often thinnest where the legal stakes are highest. Healthcare adds consent and auditability but rarely couples memory to a HIPAA-style audit cycle; finance, legal, and cybersecurity add non-repudiation and rollback obligations but score portfolio returns or verdict accuracy rather than recovery after a poisoned or lapsed-authority action. The domain view therefore reproduces the corpus-wide skew at a finer grain, naming the missing instrument per domain rather than only in aggregate.

\subsection{Web and computer-use agents}

Web and computer-use agents supply the canonical environments against which always-on claims are now measured. The foundational substrate is a set of realistic, resettable environments: WebArena provides a self-hosted web environment with reproducible sites \citep{zhou2024webarenarealisticwebenvironment}; OSWorld grounds open-ended multimodal control in a real operating system with screenshots in and keyboard or mouse out \citep{xie2024osworldbenchmarkingmultimodalagents}; AndroidWorld exposes a live Android device with programmatic tasks that read, modify, and tear down real system state \citep{rawles2024androidworld}; and BrowserGym unifies many browser-task environments into one ecosystem for cross-environment comparison \citep{lesellierdechezelles2025browsergym}. WorkArena and its enterprise extensions push the same substrate toward consequential knowledge work atop ServiceNow \citep{drouin2024workarena, boisvert2024workarena}, and TheAgentCompany frames an entire simulated software company so tasks resemble the durable, multi-day obligations of a real workplace \citep{xu2025theagentcompany}. The comprehensive survey of LLM-brained GUI agents maps this area across web, mobile, and desktop, cataloguing action spaces and memory strategies \citep{zhang2024llmguisurvey}. These environments were built to test capability, not persistence: each task resets, identity is fresh, and there is no consent boundary or recovery obligation to violate.

The persistence demand this domain adds is most visible in the systems that carry state across the resets these benchmarks impose, where two strategies have emerged. The first builds an external, retrievable record of past interaction. Agent Workflow Memory induces reusable workflow abstractions from prior trajectories so a procedure learned in one session applies in the next \citep{wang2024agentworkflowmemory}; ReasoningBank distills generalizable reasoning strategies from self-judged success and failure trajectories into a persistent, retrievable store with a closed update loop \citep{ouyang2025reasoningbank}; WebAtlas builds a persistent cognitive map of websites through exploration, storing failed interactions as cross-task memory for look-ahead planning \citep{cheng2025webatlas}; and branch-and-browse exploration shares a page-action memory within and across sessions, with web-state replay as a recovery primitive \citep{he2025branchandbrowse}. The mobile and computer-use line follows the same logic: OS-Copilot's FRIDAY self-improves by accumulating learned skills and tools and reusing them on unseen applications \citep{wu2024oscopilot}; Cradle pairs general computer control with explicit memory and skill-curation modules that persist across games and software \citep{tan2024cradle}; Agent-S introduces experience-augmented hierarchical planning that combines external knowledge search with internal experience retrieval \citep{agashe2024agents}; Mobile-Agent-E accumulates reusable Tips and executable Shortcuts across tasks \citep{wang2025mobileagente}; MobileRAG's MemRAG component retains and retrieves prior interactions so the agent stops reconstructing the interface from scratch \citep{loo2025mobilerag}; and ReinAgent organizes a multi-agent GUI system around a shared memory module with slot-based information management and proactive updates \citep{jia2025reinagent}. The second strategy folds persistence into the model's own behavior: STAMP trains agents to learn what and when to memorize and retrieve via deterministic memory variables in virtual environments \citep{wang2026stamp}, and the dual-memory M2 design for long-horizon web agents combines training-free internal trajectory summarization with an external offline insight bank \citep{yan2026m2dualmemory}. PerPilot adds a personalization layer, pairing memory-based retrieval with reasoning-based exploration to resolve under-specified personalized instructions \citep{wang2025perpilot}.

A recent result shows that this accumulation is not self-justifying. Under a token-matched budget, a vanilla baseline matches or beats online memory, skill, and workflow modules on WebArena once the retrieval cost of those modules is counted against them \citep{hajimiri2026budgetconstrained}. This cost-of-memory critique shows the value of persistent web state is conditional on governance, not volume, a domain-specific version of the corpus-wide pattern. Two further results sharpen where the governance has to live. iOSWorld is the first native iOS simulator benchmark built around a persistent user identity across twenty-six apps, with a large fraction of its tasks explicitly requiring that identity to be carried correctly \citep{jang2026iosworld}, making identity and scope first-class rather than reset away. The naive-visual-memory study shows experiential full-image visual memory helps GUI state recognition but actively hurts action grounding, isolating a perception-reasoning-action failure taxonomy on OSWorld \citep{choi2026naivevisualmemory}: accumulated state is net-negative at the act boundary, a direct instance of the actionability axis going wrong. The Darwin roadmap for self-evolving GUI agents names memory management as one of three pillars alongside task curricula and outcome verification \citep{beechey2026darwin}, treating persistence as a design axis. The problem is concrete: these environments read, modify, and tear down real device and web state, with real side effects (a sent message, a deleted record, a committed purchase), yet the benchmarks we examined do not score whether the agent can roll that side effect back when the action turns out to have been licensed by stale or poisoned state. AndroidWorld and OSWorld check final state for success, not recoverability. The web domain has a strong substrate for testing actionability and a much thinner one for testing rollback traceability.

\subsection{Embodied and robotics agents}

Embodied agents add a demand the web domain can avoid: the agent's state is a claim about a physical world that is only partially observable, drifts on its own, and cannot be reset by reloading a page. Persistence here is more than retrievable history. It is a maintained belief about object locations, world rules, and irreversible physical actions, and the benchmark literature has moved toward measuring this belief rather than task success alone. WorldLines builds a long-horizon household benchmark over action-native state trails with paired Memory-QA and Task-Planning tracks, isolating overwritten-world-state and partial-observability failures and supplying a visibility-aware memory baseline \citep{zhang2026worldlines}. Ego2World compiles egocentric cooking videos into hidden rule-governed worlds where the agent plans over a partial belief graph, isolating belief-state maintenance and showing that action-overlap metrics overestimate physical-state success \citep{cheng2026ego2world}. LMEE scores the process of memory utilization as well as the outcome for embodied exploration \citep{wang2026lmee}. These three mark a maturation: the field now grants that an embodied agent can reach a goal with a corrupted belief, a failure even when the task score is positive, an actionability-versus-provenance distinction the early navigation benchmarks could not express.

The mechanism literature responds with three contrasting approaches to keeping embodied belief trustworthy, and the contrast exposes a governance frontier. The first treats spatial memory as a high-fidelity reconstruction supporting re-observation: GS-Mem uses 3D Gaussian Splatting as persistent spatial memory to grant post-hoc re-observability so observations missed at capture time can be recovered later \citep{lu2026gsmem}, a substrate-level recoverability mechanism. The second injects calibrated uncertainty into stored semantics: the remember-with-confidence approach equips spatio-temporal memory with per-object semantic uncertainty derived from cross-view caption scatter and performs budget-constrained active refinement, so unreliable stored descriptions are flagged and resolved rather than trusted as oracle \citep{zhang2026rememberconfidence}, one of the few works in any domain that builds validation directly into the stored state. The third draws on cognitive psychology to diagnose how embodied memory fails: EMEM organizes an embodied-memory benchmark over ProcTHOR around eight cognitive paradigms (false-memory lures, pattern separation, source monitoring, serial position, long-horizon interference), separating false-memory and consolidation failures from spatial-storage failures \citep{rasheed2026emem}. The social and lifelong line extends persistence across community interaction: ELLA equips an embodied social agent with name-centric semantic plus spatiotemporal episodic memory accumulated over lifelong community life \citep{zhang2025ella}; RoboOS-Next proposes a unified memory representation for lifelong adaptation and scalable coordination across heterogeneous robot teams \citep{tan2025roboosnext}; MindForge frames lifelong cultural learning through structured theory-of-mind shared across agents \citep{lica2024mindforge}; and GOAT-Bench tests navigation to a sequence of multimodal goals within an episode, reusing memory of previously found objects \citep{khanna2024goatbench}. The cross-session stress test is RoboMME, which pads a robot's query history with unrelated sessions and shows memory-augmented vision-language-action policies decay as distractor sessions accumulate \citep{rathi2026robomme}, the embodied face of the interference failure mode.

The persistence demand here is that physical state has no undo. A web purchase can sometimes be refunded; a manipulation that knocks over a glass cannot be rolled back, only mitigated. This makes recoverability and provenance preservation the binding constraints, yet the benchmark designs above still score belief fidelity and task success, not whether the agent records enough provenance to explain a wrong physical action or maintains the recovery path that mitigation requires. EMEM diagnoses where false memories enter but does not score deletion propagation when a hallucinated object must be purged from a downstream plan; remember-with-confidence flags unreliable descriptions but does not test rollback of an action already taken on a description later found wrong. The open problem for embodied agents is therefore a governance loop matched to physical irreversibility: a write-validate step that gates which beliefs may authorize a non-reversible motor command, and an audit trace sufficient to attribute a physical error to the stored belief that caused it. No embodied benchmark in the corpus scores this.

\subsection{Healthcare}

Healthcare turns the governance gap into a legal and safety obligation. The persistence demand the domain adds is twofold: consent and deletion under regimes such as HIPAA and GDPR, and auditability of why a clinical decision was made and on what evidence. State here is the longitudinal patient record, and it carries an authority structure absent from web or robotics: the authoritative source (the chart, the FHIR record) outranks the conversational source (what the patient recalls), and conflating them is a documented safety hazard. The benchmark family has converged on this longitudinal, multi-stage character. AgentClinic was the seminal clinical-agent benchmark to study a persistence axis explicitly, with a notebook memory that survives across cases \citep{schmidgall2024agentclinic}. MedAgentBench defines the EHR-agent benchmark family with an actionability and scope axis built in, since its tool calls write irreversibly against patient records \citep{jiang2025medagentbench}. ClinEnv benchmarks an LLM as attending physician over an ordered sequence of irreversible decision stages per admission, making sequential irreversibility first-class \citep{lu2026clinenv}. MediLongChat tests recall and reasoning over a patient's longitudinal history with a dedicated cross-dialogue reasoning task \citep{hu2026medilongchat}, and ESMemEval defines a longitudinal user-memory benchmark family with explicit state axes spanning information extraction and temporal reasoning \citep{chen2026esmemeval}. PsychEval frames counseling as a longitudinal task requiring sustained memory and dynamic goal tracking across sessions \citep{pan2026psycheval}, a framing TheraMind operationalizes with an explicit cross-session persistent-state lifecycle: write after a session, retrieve within a session, and an update strategy across the therapeutic relationship \citep{hu2025theramind}.

The clinical mechanism literature contains some of the most governance-relevant ideas in any application domain, because clinicians will not accept an agent that silently overwrites a record. MedMemoryBench directly targets always-on memory in clinical agents and formalizes the memory-saturation failure mode, where sustained information inflow degrades the agent's ability to use its own memory \citep{wang2026medmemorybench}, a continual-operation degradation result with clinical stakes. VitalTrace replaces unbounded patient-history serialization in an ICU multi-agent framework with a compact persistent patient-state memory whose updates obey curated physiological state-transition rules, improving temporal consistency over long trajectories \citep{qu2026vitaltrace}, a rare instance of a validation step (rule-governed transitions) gating writes. TrajOnco performs temporal reasoning over sequential clinical events through a multi-agent long-term-memory framework \citep{zeng2026trajonco}. The strongest governance pattern is the dual-stream clinical memory that separates the patient's recall-biased self-report from the authoritative FHIR record and runs a reconciliation engine to flag discrepancies, rejecting the overwrite-with-latest-statement default as a safety hazard \citep{pugh2026dualstream}. This is authority monotonicity and provenance preservation realized in a deployed-style clinical design: the authoritative source is never silently superseded by a lower-authority statement, and discrepancies are surfaced rather than resolved by recency.

Yet even here the gap persists: these systems build provenance separation and rule-gated updates, but none couples the memory lifecycle to the audit and deletion cycle the regulation actually requires. A right-to-be-forgotten request under GDPR, or a record-amendment obligation under HIPAA, demands that a deletion propagate through every derived tier (summaries, embeddings, cached prompts, inter-agent messages, fine-tuned adapters) and that the propagation be verifiable; no clinical benchmark in the corpus scores deletion completeness across these tiers, and no clinical system demonstrates rollback of a clinical action after the authorizing record is amended. TheraMind maintains a therapeutic relationship across sessions but offers no erasure guarantee; ClinEnv stages irreversible decisions but does not test recovery after a decision is later found to rest on an amended record. Domain-specific governance, HIPAA-style audit trails coupled to the memory lifecycle and GDPR-style verifiable erasure, is the largest unaddressed demand here, which is why the dual-stream and rule-gated systems read as exceptional rather than standard.

\subsection{Finance, legal, and cybersecurity}

Finance, legal, and cybersecurity agents share a persistence demand that healthcare only partially carries: non-repudiation. A trade, a filing, or a credential grant is an act whose record must be tamper-evident and attributable after the fact, and whose reversal (where reversal is even possible) must be traceable. The state these agents hold is a ledger of consequential commitments, and the binding axes are authority (who licensed the action), provenance (what evidence and which credentials backed it), and rollback traceability (can the action and its derived state be unwound, and is the unwinding itself recorded). The finance literature is the oldest in this group. FinMem is the canonical trading agent whose core contribution is a layered deep, intermediate, and shallow memory module aligned to decision timescales \citep{yu2023finmem}; FinAgent extends this with a diversified memory-retrieval system over historical market data plus a dual reflection module \citep{zhang2024finagent}; FinCon is a manager-analyst multi-agent system whose risk-control component episodically self-critiques to update persistent investment beliefs \citep{yu2024fincon}; and InvestorBench was the first standardized benchmark for memory-equipped financial decision agents across products and market environments \citep{li2024investorbench}. Two recent benchmarks reframe finance around state explicitly: $\tau$-Banking extends $\tau$-bench into a fintech customer-support domain where agents must navigate roughly seven hundred policy and procedure constraints \citep{shi2026tauknowledge}, and the thousand-day-scale RetailBench simulates a single store's pricing, replenishment, inventory aging, and cash flow, where errors compound over the horizon \citep{zhang2026retailbench}. KTD-Fin sharpens a provenance subtlety the other finance work elides, separating pretraining-memorized market data, fixed at the knowledge cutoff, from genuinely retrieved current state \citep{zhu2026ktdfin}, a provenance-axis question (is this belief sourced from current evidence or a stale parametric prior) that directly governs whether a trade is justified.

The legal line is younger but unusually aligned with the persistent-state frame because legal procedure is, by construction, a governed state machine. LegalWorld models Chinese civil litigation as a causally connected five-stage state chain built on tens of thousands of paired judgments \citep{zuo2026legalworld}; SimCourt was the first framework to model the procedural structure of Chinese criminal trials with five courtroom roles whose legal agents carry role-bound state \citep{zhang2025simcourt}; and the civil-court multi-agent simulation integrates a memory module with statute retrieval to support long-process adjudication \citep{chen2026civilcourt}. These systems implicitly demand authority monotonicity (a later stage cannot expand the authority granted by an earlier ruling) and provenance preservation (a verdict must be traceable to the evidence and statutes that produced it), because that is how legal procedure works; but they are evaluated on adjudication accuracy, not on whether the agent's state respects those procedural invariants under perturbation. Table~\ref{tab:resources-finlegal} contrasts what each representative resource measures against the governance obligation the domain imposes.

\begin{table}[t]
\centering
\small
\setlength{\tabcolsep}{4pt}\renewcommand{\arraystretch}{1.2}
\begin{tabularx}{\linewidth}{p{0.17\linewidth}YYY}
\toprule
Domain & Representative resource & Persistence demand added & Governance gap \\
\midrule
Finance (trading) & FinMem \citep{yu2023finmem}, FinAgent \citep{zhang2024finagent}, InvestorBench \citep{li2024investorbench} & Non-repudiation of trades; timescale-aligned belief & Scores returns, not rollback or audit of a wrong trade \\
Finance (current vs. prior) & KTD-Fin \citep{zhu2026ktdfin} & Provenance of belief: current evidence vs. parametric prior & No deletion or correction of a stale-prior-driven action \\
Finance (long horizon) & RetailBench \citep{zhang2026retailbench} & Compounding over thousand-day operations & No recovery instrument for compounded state error \\
Legal procedure & LegalWorld \citep{zuo2026legalworld}, SimCourt \citep{zhang2025simcourt} & Procedural state chain; role-bound authority & Scores verdict accuracy, not invariant preservation under perturbation \\
Cybersecurity-adjacent & $\tau$-Banking \citep{shi2026tauknowledge} & Policy-constrained action over hundreds of rules & No test of credential rollback after lapsed authority \\
\bottomrule
\end{tabularx}
\caption{Finance, legal, and security resources read through the non-repudiation and rollback demand each adds. The recurring gap is that the benchmarks score the outcome (return, verdict, policy compliance) but not the return arc of the lifecycle (audit, deletion, rollback) that non-repudiation requires.}
\label{tab:resources-finlegal}
\end{table}

The gap statement for this group is the most direct in the part. Non-repudiation is, in lifecycle terms, the conjunction of provenance preservation and rollback traceability: every consequential act must be attributable and, where reversed, the reversal must itself be auditable. The systems above build sophisticated memory for deciding what to do and essentially nothing for unwinding what was done: none scores rollback of a tool-originated commitment, a placed trade, a granted credential, a filed document, after the authority that licensed it lapses or is found invalid, the exact open instrument the broader corpus identifies, a retried or reversed action must not resurrect already-consumed authority or duplicate an irreversible side effect. Where the side effect is money, a legal filing, or a security credential, the absence of this instrument is not an academic gap but a deployment blocker, unaddressed across every resource in this group.

\subsection{Education and tutoring}

Education agents add a persistence demand that looks gentle but is governance-heavy: the authoritative state is a model of what the learner knows, and that model must mutate correctly over time. A tutoring agent that treats a once-mastered skill as permanently mastered, or fails to decay knowledge the student has since forgotten, will mis-instruct; an agent that overwrites a misconception rather than verifying it has been corrected will record success where there is none. The binding axes are mutability (the learner state evolves at the speed of learning and forgetting) and provenance (the basis for believing a skill is mastered), and the mechanism literature is unusually explicit about modeling this evolving state. TutorLLM provides a concrete persistent learner-state substrate, using a knowledge-tracing model to predict per-student knowledge state as the memory the tutor conditions on \citep{li2025tutorllm}. TASA couples an event memory and persona profile with a continuous forgetting curve over knowledge tracing, deliberately decaying previously-mastered skills so instruction is recalibrated to the learner's current mastery \citep{wu2025tasa}, one of the clearest domain instances of principled, relevance-driven forgetting being a feature rather than a failure. DeepTutor builds an agentic tutor with a dynamic learner memory that adapts to evolving needs and ships TutorBench with learner profiles \citep{zhao2026deeptutor}; AgentTutor formalizes a persistent learner-state lifecycle with organize and retrieve operations over a learner profile plus experience memory \citep{liu2026agenttutor}; and PsychAgent combines persistent episodic and experiential memory with weight-level consolidation, internalizing skills via rejection fine-tuning \citep{yang2026psychagent}, a parametric-plus-non-parametric hybrid that raises the provenance question of where a consolidated skill came from.

Contrasting these makes a mechanism point. TASA and TutorLLM treat learner state as something tracked and decayed by a principled model, so mutation is governed by a knowledge-tracing prior; DeepTutor and AgentTutor treat it as adaptive memory shaped by interaction, so mutation is governed by what the agent observes. The former is more auditable (decay follows a curve), the latter more flexible (it captures idiosyncratic learner change), and neither validates the key correctness condition: that a recorded mastery update reflects genuine understanding rather than a lucky answer, and that a recorded misconception was corrected rather than overwritten. The gap is validation at the write boundary and recoverability of the learner model. Education has, in TASA's forgetting curve, one of the better forgetting mechanisms in the corpus, but no education benchmark scores whether an erroneous mastery update can be detected and rolled back, and FERPA-style governance over the durable learner record, who may read, amend, or erase it, is absent here. The learner state is the most sensitive long-horizon record after the medical chart, and it is governed least.

\subsection{Coding agents}

Coding agents make the persistence demand concrete and testable in a way few other domains do, because a codebase is a large, structured, versioned state with an existing notion of rollback (the commit graph) that the agent's memory must respect. The demand the domain adds is cross-session coherence: an agent must hold a project's conventions, prior fixes, and codebase structure across sessions without repeating known mistakes or violating learned conventions. The codified-context work documents the core persistence failure directly: coding agents lose coherence and conventions across sessions and repeat known mistakes \citep{vasilopoulos2026codifiedcontext}. The systems answer with hierarchical persistent memory. Prometheus builds persistent codebase-navigation memory beyond the context window, exercising the organize and retrieve lifecycle and the scope axis over a repository \citep{pan2025prometheus}; MemRepair maintains a hierarchical persistent memory that reuses past fixes and learns from failed validation feedback across repair attempts \citep{liu2026memrepair}, one of the few systems anywhere that treats a failed validation as a first-class write rather than a discarded outcome. On the benchmark side, coding contributes useful always-on stress tests. MemoryCode is a synthetic multi-session coding benchmark that tests whether an agent can retrieve and act on instructions across sessions amid distractors \citep{rakotonirina2025memorycode}, and SWE-Bench-CL restructures SWE-Bench-Verified GitHub issues into chronological per-repository task streams, pairing a semantic-memory agent with forgetting and forward and backward-transfer metrics and a composite continual-learning score \citep{joshi2025swebenchcl}. SWE-Bench-CL is notable because it imports continual-learning evaluation, transfer and forgetting metrics, into a domain with real ground truth (does the patch pass the test), making it one of the rare resources that measures whether accumulation compounds or interferes rather than assuming it helps.

The contrast with the web domain is informative: coding agents operate over state that already has version control, so the rollback primitive exists at the substrate, yet the memory layer does not use it. MemRepair learns from failed validations but does not roll its own memory back to a known-good state when a learned fix later proves wrong; Prometheus navigates and reasons but does not check consistency when it modifies shared code state. The open problem is to bind the agent's persistent memory to the codebase's own rollback semantics, so reverting a commit also reverts the memory writes derived from it (deletion propagation), and a poisoned procedural memory, a learned anti-pattern presented as a fix, can be quarantined and traced (provenance preservation, rollback traceability). The domain has the cleanest available substrate for studying lifecycle-complete memory governance and has not yet used it for that purpose.

\subsection{Ambient, wearable, and IoT agents}

Ambient, wearable, and IoT agents most fully realize the always-on premise, because the agent literally never resets: it observes a continuous, multimodal life stream and is expected to surface help proactively. The persistence demands it adds are three. First, capture is unbounded and continuous, so the write and consolidation boundaries face the heaviest pressure in any domain. Second, the agent must decide when to act unprompted, making trigger state and proactivity calibration a first-class state type rather than an afterthought. Third, the stream is intimate first-person life data, so privacy and scope are acute. The benchmark literature here is the newest and most directly aimed at the always-on requirement. EgoLife provides a three-hundred-hour, six-participant, one-week egocentric AI-glasses dataset and the EgoLifeQA benchmark for a life assistant that recalls past events and monitors habits, with an EgoButler retrieval-memory system \citep{yang2025egolife}. TeleEgo benchmarks always-on egocentric assistants over fourteen-plus hours of synchronized streaming multimodal data and, distinctively, scores Memory Persistence Time and Real-Time Accuracy across twelve subtasks \citep{yan2025teleego}, naming time-to-persistence as a metric rather than only recall. LifeDialBench pushes to a one-year simulated horizon plus a seven-day real EgoMem, evaluated under an online temporal-causality protocol, and reports a finding that anchors this domain discussion: lossy sophisticated memory systems underperform a high-fidelity RAG baseline \citep{zheng2026lifedial}. This is the ambient-domain echo of the cost-of-memory critique: elaborate consolidation can be net-negative when it drops the handles a later query needs.

The mechanism and proactivity line is where this domain contributes a state type the others underweight. ContextAgent fuses open-world sensory perception with personas built from historical data to predict when proactive service is needed and then calls tools, releasing the ContextAgentBench benchmark \citep{yang2025contextagent}; Memento is a wearable AR assistant that permanently captures verbal queries with spatiotemporal and activity context and proactively re-surfaces recurring interests when a matching context recurs \citep{kim2026memento}; and EgoSelf builds a graph-based interaction memory of a user's first-person history to capture long-term habits and predict future interactions \citep{wang2026egoself}. The smart-home and IoT line extends the same proactivity into device control under constraints: PersonalHomeBench evaluates foundation-model agents in personalized homes by progressively building evolving household state and scoring both reactive and proactive ability \citep{bharadwaj2026personalhomebench}; IoTGPT memorizes decomposed instruction subtasks for reuse to cut LLM calls while adapting them to per-user preferences \citep{yu2026iotgpt}; AirAgent maintains a dynamic memory-tag layer that continuously updates a personalized user profile and fuses it with live sensor data to drive proactive device control under explicit constraints \citep{men2026airagent}; and the shopping-companion benchmark, though e-commerce, shares the ambient pattern of persisting cross-session preference over a 1.2M-item pool and isolating preference-hallucination cascades unique to long-horizon purchasing \citep{yu2026shoppingcompanion}.

The governance gap is acute for ambient agents because the demands are high and the instruments remain young. Proactivity calibration, the trigger-state question of when to surface help versus stay silent, is named (ContextAgent, the broader proactive-agent line \citep{lu2024proactiveagent}) but not governed: no ambient benchmark we examined scores the authority and scope conditions under which a proactive action is licensed, so an over-eager surfacing (intrusion) and a missed-help failure are treated as accuracy rather than scope or authority violations. Privacy is acute and under-addressed: the life stream is the most sensitive durable record in any domain, yet ambient resources rarely score deletion completeness across the derived tiers a wearable accumulates (raw video, transcripts, summaries, habit graphs, proactive-trigger indices), and a right-to-be-forgotten request against a year of lifelog has no demonstrated propagation path. LifeDialBench's finding that lossy consolidation underperforms high-fidelity retrieval is the consolidation-boundary failure made domain-specific: always-on memory needs consolidation that remains addressable by retrieval, with the handles preserved, and the ambient domain is where the pressure to drop them is highest and the cost best measured.

\subsection{Scientific discovery agents}

Scientific-discovery agents add a persistence demand distinct from the others: the durable state is the research program itself, the accumulating record of what has been tried, what worked, and what failed, and its governance question is whether that record steers exploration correctly over a campaign rather than misleading it. Two representative systems make the structure explicit. DeepScientist's main state object is a cumulative, tiered Findings Memory in which validated findings accumulate, are promoted to higher-fidelity validation, and steer a Bayesian explore-exploit policy across a month-long discovery campaign \citep{weng2025deepscientist}. EvoScientist maintains separate persistent ideation and experimentation memories that record both feasible and previously-failed directions so research strategies improve across the program \citep{lyu2026evoscientist}. Both treat the negative result, the failed direction, as a first-class durable write, a provenance-rich design: the value of the memory depends on faithfully attributing each finding to the evidence and validation tier that produced it. DeepScientist's promotion mechanism is, in lifecycle terms, a validation-gated update: a finding does not gain authority over the policy until it survives higher-fidelity validation, an authority-monotonicity discipline applied to scientific belief.

The gap is recoverability of the program when a promoted finding turns out to be wrong. A finding that was validated, promoted, and then used to steer dozens of subsequent experiments contaminates everything downstream of it; rolling it back requires deletion propagation through the explore-exploit policy and every derived finding, and an audit trace that can attribute the wasted exploration to the retracted result. Neither system scores this recovery, nor demonstrates that a retraction propagates correctly, which matters because a cumulative findings memory works by letting early findings compound into later ones. The same path also lets a single contaminated finding do maximum damage. Scientific discovery therefore repeats the corpus-wide finding under unusually clear conditions: the systems are strong at accumulating and steering, and untested at recovering when the accumulation is poisoned.

\subsection{Memory frameworks, libraries, and serving substrates as reusable resources}

Cutting across every domain above is a layer of reusable resources, the memory frameworks, libraries, and serving substrates that domain systems build on, and the way the domains consume this layer is itself diagnostic. The dominant pattern is that domain systems adopt a retrieval-and-storage library (a vector store, a temporal knowledge graph, an extracted-fact memory) and inherit its governance properties wholesale, which usually means they inherit little governance by default. Supporting evidence comes from the budget-matched web study and the ambient-domain finding together: under a token-matched budget the online memory, skill, and workflow modules do not beat a vanilla baseline on WebArena \citep{hajimiri2026budgetconstrained}, and at the lifelog horizon a lossy sophisticated memory system underperforms a high-fidelity RAG baseline \citep{zheng2026lifedial}. As a statement about the resource layer, these suggest that popular memory libraries can add cost and lossiness without the governance that would justify them: they implement write, organize, and retrieve and omit validate, audit, deletion-propagation, and rollback. A domain builder who reaches for one gets accumulation and retrieval and must build the return arc themselves, which is rarely visible in the domain literature above.

The serving and substrate question, where state physically lives and how it is bounded under load, is also consumed cross-domain rather than solved per-domain, and it is the layer where the continual-operation degradation modes surface. MedMemoryBench's memory-saturation finding in healthcare \citep{wang2026medmemorybench} and RoboMME's distractor-session decay in robotics \citep{rathi2026robomme} are the same phenomenon, unbounded growth degrading retrieval, observed through different domains, which argues that bounded, evictable, governed serving substrates are a shared resource need rather than a per-domain one. The few systems that do treat the substrate as governable (VitalTrace's rule-governed compact patient state \citep{qu2026vitaltrace}, GS-Mem's re-observable spatial substrate \citep{lu2026gsmem}) are domain-specific and not packaged as reusable resources, so the next clinical or embodied builder cannot reuse them. The gap at this layer is therefore an interoperability and governance gap at once: the works we examined do not yet provide a reusable memory substrate that ships the return arc, validation at write, provenance through consolidation, verifiable deletion propagation, and rollback traceability, as first-class operations, so a domain system inherits governance less readily than it inherits retrieval. Until such a substrate exists, each domain tends to re-derive accumulation cheaply while leaving governance bespoke. The resource layer turns the survey's diagnosis into an engineering target: governed persistent state should be a reusable substrate, not a per-paper afterthought.
